% ============================================================
% Plantilla: Artículo IEEE para conferencias (clase IEEEtran, dos columnas)
% Compila con pdflatex + bibtex.
% ============================================================
\documentclass[conference]{IEEEtran}

\usepackage[utf8]{inputenc}
\usepackage[T1]{fontenc}
\usepackage[spanish,es-tabla]{babel}
\usepackage{amsmath,amssymb}
\usepackage{graphicx}
\usepackage{booktabs}

% IEEEtran no traduce este título con babel
\renewcommand{\IEEEkeywordsname}{Palabras clave}

\begin{document}

\title{Título del artículo para la conferencia}

\author{
  \IEEEauthorblockN{Primer Autor}
  \IEEEauthorblockA{\textit{Departamento} \\
    \textit{Universidad}\\
    Ciudad, País \\
    correo@ejemplo.com}
  \and
  \IEEEauthorblockN{Segundo Autor}
  \IEEEauthorblockA{\textit{Departamento} \\
    \textit{Universidad}\\
    Ciudad, País \\
    correo@ejemplo.com}
}

\maketitle

\begin{abstract}
Resume el problema, el método, los resultados y la conclusión en un solo párrafo
de no más de 250 palabras.
\end{abstract}

\begin{IEEEkeywords}
palabra uno, palabra dos, palabra tres
\end{IEEEkeywords}

\section{Introducción}
Presenta el contexto y la contribución principal del trabajo \cite{lamport1994}.

\section{Trabajos relacionados}
Resume los enfoques existentes y en qué se diferencia tu propuesta \cite{knuth1984}.

\section{Método propuesto}
Describe la propuesta. Las ecuaciones se numeran automáticamente:
\begin{equation}
  f(x) = \sum_{k=0}^{n} a_k x^k
  \label{eq:polinomio}
\end{equation}

\section{Resultados experimentales}
La Tabla~\ref{tab:resultados} resume los resultados obtenidos con~\eqref{eq:polinomio}.

\begin{table}[htbp]
  \caption{Resultados experimentales}
  \label{tab:resultados}
  \centering
  \begin{tabular}{lcc}
    \toprule
    Configuración & Precisión & Recall \\
    \midrule
    Base & 0.81 & 0.77 \\
    Propuesta & \textbf{0.89} & \textbf{0.85} \\
    \bottomrule
  \end{tabular}
\end{table}

\section{Conclusiones}
Resume los aportes y el trabajo futuro.

\bibliographystyle{IEEEtran}
\bibliography{referencias}

\end{document}
